\documentclass[10pt,a4paper]{article}
\usepackage[left=2.5cm,right=2.5cm,top=2.5cm,bottom=2.5cm,]{geometry}
\usepackage{amsmath,amssymb,amsthm}
\newtheorem{theorem}{Theorem}[section]
\newtheorem{lemma}[theorem]{Lemma}
\newtheorem{proposition}[theorem]{Proposition}
\theoremstyle{definition}
\newtheorem{definition}[theorem]{Definition}
\usepackage{setspace}
\setlength{\parindent}{0pt}
\onehalfspacing
\begin{document}
\begin{definition}
    $G$ is a group, we 
    call $H$ is the subgroup of $G$ if and only if $H\subseteq G\land\forall x\forall y(x\in H\land y\in H\rightarrow x^{-1}\in H\land xy\in H)$
\end{definition}
Obviously, we can see $|H|\le|G|$. Another useful proposition is:
\begin{proposition}
    $H\subseteq G$, if $H$ is the subgroup of $G$ if and only if $\exists x(x\in H)\land\forall x\forall y(x\in H\land y\in H\rightarrow xy^{-1}\in H)$.
\end{proposition}
\begin{proof}
    $H$ is the subgroup of $G$, for $H$, the operation is closed obviously.
    let $y=x^{-1}$, $xy=xx^{-1}=\mathrm{id.}\in H$. $\forall x\in H,\mathrm{id.}\in H,\mathrm{id.}a^{-1}=a^{-1}\in H$.
    $\forall x,y\in H$, $a(b^{-1})^{-1}=ab\in H$. $H\le G$.
\end{proof}
In fact, this can itself be viewed as a right regular action. Consider $\pi:G\times G\to G$, $(g,x)\mapsto g\cdot x:=xg^{-1}$.
associative law $\forall g,h\in G$, $(gh)\cdot x=xh^{-1}g^{-1}$, $g\cdot(h\cdot x)=g\cdot xh^{-1}=xh^{-1}g^{-1}$. 
This tells us $H$ is the subgroup of $G$ iff. $H$ is a fixed subset of $G$ under self action. Next, I will offer another perspective on subgroups.
Consider retract map $\pi:G\to H$, $\pi\mid_H=\mathrm{id_H.}$. $\forall h(h\in H\rightarrow\pi(h)=h)$. let $N=\ker\pi$, obviously, $N\cap H=\{\mathrm{id.}\}$.
$N\trianglelefteq G$, so the elements of $G$ can be uniquely represented as $\forall g\in G$, $g=(g\cdot(\pi(g))^{-1})\cdot\pi(g)$.
so we have \[1\longrightarrow N\overset{\iota}\longrightarrow\ G\overset{\pi}\longrightarrow H\longrightarrow1\]
$G\cong N\rtimes H$.
\begin{proposition}
    $A$, $B$ are the subgroups of $G$, if $A\cup B$ is the subgroup of $G$ iff. $A\subseteq B\lor B\subseteq A$.
\end{proposition}
\begin{proof}
    if $A\cap B=\{\mathrm{id.}\}$, the elements in $A$ and $B$ are closed under their respective operations. Obviously, $A\cup B$ is not subgroup.
    so $A\cap B\ne\{\mathrm{id.}\}$. if $A\subseteq B$, $\forall x\in B$, $x\in A$, Obviously.
\end{proof}
\begin{definition}
    Let $G$ is an abelian group, the torsion group $T$ is the subgroup of $G$, where $T=\{g\in G\mid|g|<\infty\}$.
\end{definition}
Easily, we have
\[0\longrightarrow T(G)\overset{\iota}\longrightarrow G\overset{\pi}\longrightarrow G/T(G)\longrightarrow0\]
if $G$ is finite, then $T(G)=G$, $G/T(G)=\{0\}$, $G\cong G\oplus\{0\}$.
if $G$ is infinite, $G\cong T(G)\oplus G/T(G)$.(It not always correct.)
This need more powerful condition, specficty, $\exists\varphi:G/T(G)\to G$, $\pi\circ\varphi=\mathrm{id_H.}$.
\begin{proposition}
    $H$ is the subgroup of $G$, $H\le N_G(H)$.
\end{proposition}
\begin{proof}
    $H\curvearrowright H$,$N_H(H)$,$\forall h,x\in H,x\mapsto hxh^{-1}$, $hxh^{-1}\in H$, $hHh^{-1}=H$, $H=N_{H}(H)$, $H=N_H(H)\le N_G(H)$.
\end{proof}
\begin{proposition}
    $A\subseteq B\subseteq G$, $C_G(B)\subseteq C_G(A)$.
\end{proposition}
\begin{proof}
    $\forall x\in G$, if an element can commute with $a\forall a\in A$, it must commute with $b\forall b\in B$.
    if an element can commute with $b\forall b\in B$, it may commute with $a\forall a\in A$. $C_G(B)\subseteq C_G(A)$.
\end{proof}
more generally, given a group $G$, some set $A\subseteq B\subseteq C\subseteq\cdots\subseteq G$. $C_G(G)\subseteq\cdots\subseteq C_G(B)\subseteq C_G(A)$.
\begin{proposition}
    $A\subseteq G$, $H\le G$, $N_H(A)=N_G(A)\cap H$.
\end{proposition}
\begin{proof}
    $\forall x\in N_H(A)$,$x\in H\land xHx^{-1}=H$,$x\in N_G(A)\land x\in H$,$x\in N_G(A)\cap H$. $x\in N_G(A)\cap H$,$x\in H\land x\in N_G(A)$. $xAx^{-1}=A$,$x\in N_H(A)$, $N_H(A)=N_G(A)\cap H$
\end{proof}
\begin{proposition}
    $A\subseteq G$,$Z(G)\le N_G(A)$.
\end{proposition}
\begin{proof}
    $\forall g\in Z(G)$,$\forall a\in A$, $gag^{-1}=a$, $gAg^{-1}=A$, $g\in N_G(A)$. $Z(G)\subseteq N_G(A)$. $Z(G)\le G\land N_G(A)\le G$,$Z(G)\le N_G(A)$.
\end{proof}
\begin{definition}
    A group $H$ is called cyclic if it can be generated by a single element. i.e. $H=\{x^n\mid n\in\mathbb{Z}\}$. 
\end{definition}
It's easy to find $H=\{x^n\mid n\in\mathbb{Z}\}=\{x^{-n}\mid n\in\mathbb{Z}\}$. That's beacuse $x^{-n}=(x^{-1})^n$.
Obviously, if $H=\langle x\rangle$, $|H|=|x|$. More specifically, if $|H|<\infty$, $\forall x\in H$ is distinct. suppose $x^a=x^b$, $a<b$,
$x^{b-a}=1$. if $|H|=\infty$, $\forall n\ne0$,$x^n\ne1$,$\forall a\ne b$,$x^a\ne x^b$,$\forall a,b\in\mathbb{Z}$. For a cyclic group, it's very useful to attention its smallest respective element. We will see in next part.
\begin{proposition}
    Given a group $G$, $x\in G$, if $x^m=1,m\in\mathbb{Z}$, $|x|\mid m$. 
\end{proposition}
\begin{proof}
    let $|x|=n$, $m=nq+r$, $x^r=x^{m-nq}$, $r=0$, $m=nq$, $n\mid m$.
\end{proof}
\begin{proposition}
    if $\langle x\rangle$ and $\langle y\rangle$ are finite and have same order, then they are isomomphism.
\end{proposition}
\begin{proof}
    Consider a map $\pi:\langle x\rangle\to\langle y\rangle$, $\forall k\in\mathbb{Z}$, $x^k\mapsto y^k$. $\pi(x^k)=y^k$. First, we will check the map is well-defined.
    Consider if $\forall a,b\in\mathbb{Z}$, $x^a=x^b$, $\pi(x^a)=\pi(x^b)$, $x^{a-b}=1$, $n\mid a-b$, $\exists q\in\mathbb{Z},a=nq+b$, $\pi(x^{nq+b})=\pi(x^b)=y^{nq+b}=y^b$. It's well-defined.
    Next, we will proof the map is homomorphic. $\forall a,b\in\mathbb{Z}$, $\pi(x^ax^b)=\pi(x^{ab})=y^{ab}$. $\pi(x^a)\pi(x^b)=y^ay^b=y^{ab}$. $\pi(x^ax^b)=\pi(x^a)\pi(x^b)$. $y^k$ is the image of $x^k$, so it's surjective. Also, they have the same order so it's bejective.
\end{proof}
\begin{proposition}
    if $\langle x\rangle$ is infinite, so $\mathbb{Z}\cong\langle x\rangle$.
\end{proposition}
\begin{proof}
    Consider a map $\pi:\mathbb{Z}\to\langle x\rangle$, $\forall k\in\mathbb{Z}$, $\pi(k)=x^k$. $\forall a,b\in\mathbb{Z}$, $a\ne b$, $x^a\ne x^b$, so it's injective. By definition, it's surjective. so it's bejective.
\end{proof}
\begin{proposition}
    Let $G$ is a group, $x\in G$, $a\in\mathbb{Z}-0$, if $|x|=n<\infty$, then $|x^a|=\frac{n}{\gcd(n,a)}$
\end{proposition}
\begin{proof}
    $|x^a|=m$,$d=\gcd(n,a)$,$n=dn'$,$a=da'$,$\forall n',a'\in\mathbb{Z}$,$\gcd(n',a')=1$, $(x^a)^{n'}=x^{an'}=x^{da'n'}=1^{a'}=1$,$m\mid n'$.
    $x^{am}=x^n$,$n\mid am$,$dn'\mid da'm$,$n'\mid a'm$, $n'\mid m$
    $m=n'$.
\end{proof}
\begin{proposition}
    $H=\langle x\rangle$, if $K\le H$, $K=\{1\}\lor K=\langle x^d\rangle$, where $d$ is the smallest positive number such $x^d\in K$.
\end{proposition}
\begin{proof}
    $H=\langle x\rangle$,$K\le H$,$\exists a\ne0$,$x^a\in K$. if $a<0$, $x^{-a}=(x^{a})^{-1}\in K$.
    $\mathcal{P}=\{b\mid b\in\mathbb{Z}^+\land x^b\in K\}$. $x^d\in K$,$\langle x^d\rangle\le K$. 
    $K\le H$,$a=qd+r$,$x^a=x^{qd+r}$,$x^r=x^{a-qd}$,$a=qd$,$x^a=(x^{d})^q\in\langle x^d\rangle$. $K\le\langle x^d\rangle$. $K=\langle x^d\rangle$.
\end{proof}
\begin{proposition}
    $H=\langle x\rangle$, $|x|=n<\infty$, $H=\langle x^a\rangle$ iff. $\gcd(a,n)=1$.
\end{proposition}
\begin{proof}
    $|x|=n<\infty$,$H=\langle x^a\rangle\iff\gcd(a,n)=1$.
    $|x^a|=|x|\iff\frac{n}{\gcd(a,n)}=n\iff\gcd(a,n)=1$
\end{proof}
This is very useful, it tells us $H$ is generated by $\varphi(n)$. For example, consider a group $\langle\mathbb{Z}/n\mathbb{Z},\times\rangle$, $H=\{\varphi(n)\mid n\in\mathbb{Z}\}$.
\begin{proposition}
    $H=\langle x\rangle$, $|H|=n<\infty$, $\forall a\in\mathbb{Z},a\mid n$ is a unique subgroup of $H$ of order $a$.
\end{proposition}
\begin{proof}
    let $d=a\mid n$, $\langle x^d\rangle$ is a subgroup of $H$ of order $a$, suppose $\exists b\in\mathbb{Z}$, $\langle x^b\rangle$ is a subgroup of $H$ of order $a$.
    $\frac{n}{d}=a=|x^b|=\frac{n}{\gcd(b,n)}$, $\gcd(b,n)=d$, $d\mid b$, $x^b\in\langle x^d\rangle$. More generally, we have $\langle x^b\rangle\le\langle x^{\gcd(b,n)}\rangle$.
\end{proof}
From this, we can see the subgroup of finite cyclic group 
\end{document}


